\section{Critical-height relationships on the same rectangles}
\label{u:sec:compatible-rectangles}

A finite rectangle selects temporal and spatial orders from the datum-generated
mixed jet, retaining their velocity and pressure coordinates and governing
relationships.  Across all finite \(N,M\), every larger rectangle restricts to
the same derivatives and pressure coordinates in each smaller rectangle, and
the family exhausts every finite mixed index.  Independently, the completed
critical-height identities are evaluated on those same finite rectangles.
Exhausting their heights gives every spatial Sobolev row, and the
Navier--Stokes recurrence reconstructs the same temporal and pressure rows.


\begin{proposition}[Kinetic-energy identity]
\label{u:prop:kinetic-energy}
For every \(0\le t<T_*\),
\begin{equation}
 \frac12\norm{u(t)}_2^2
 +\nu\int_0^t\norm{\nabla u(s)}_2^2\dd s
 =\frac12\norm{u_0}_2^2.
 \label{u:eq:kinetic-energy}
\end{equation}
Consequently, for \(0<T<T_*\),
\begin{align}
 \norm u_{L^\infty(0,T;L^2)}&\le\norm{u_0}_2,
 \label{u:eq:energy-Linf}\\
 \norm u_{L^2(0,T;H^1)}^2
 &\le\left(T+\frac1{2\nu}\right)\norm{u_0}_2^2.
 \label{u:eq:energy-L2H1}
\end{align}
\end{proposition}

\begin{proof}
Take the \(L^2\) pairing of the momentum equation with \(u\).  Integration
by parts and \(\nabla\cdot u=0\) give
\[
 \ip{(u\cdot\nabla)u}{u}_{L^2}=0,
 \qquad
 \ip{\nabla p}{u}_{L^2}=0,
 \qquad
 \ip{\Delta u}{u}_{L^2}=-\norm{\nabla u}_2^2.
\]
Integration in time proves \eqref{u:eq:kinetic-energy}.  The two estimates
follow by discarding the nonnegative dissipation and adding
\(\int_0^T\norm{u(s)}_2^2\dd s\).
\end{proof}

Pairing a differentiated velocity equation with a spatial derivative of
that row expresses the change of its Sobolev energy through viscous
dissipation and the full transport--pressure contribution.  Repeated time
differentiation then exposes higher spatial energies, with the differentiated
nonlinear terms retained in each identity.

\begin{definition}[Completed height rows]
\label{u:def:critical-height-rows}
For \(r\in\Nzero\) and \(s\ge0\), set
\begin{align}
 E_{r,s}(t)&:=\frac12\norm{\Lambda^sV_r(t)}_2^2,
 \label{u:eq:temporal-row-spatial-energy}\\
 \cP_{r,s}(t)&:=-\operatorname{Re}
 \ip{\Lambda^sV_r(t)}
 {\Lambda^s\!\left(
   \sum_{a=0}^{r}\binom ra(V_a\cdot\nabla)V_{r-a}+\nabla Q_r
 \right)(t)}_{L^2}.
 \label{u:eq:temporal-row-pressure-complete-transfer}
\end{align}
For the base row write \(E_s:=E_{0,s}\), \(\cP_s:=\cP_{0,s}\), and
\(H:=E_{1/2}\).
\end{definition}

Differentiating this quadratic height along the classical history gives
\[
 H^{(r)}(t)=\frac12\sum_{a=0}^r\binom ra
 \ip{\Lambda^{1/2}V_a(t)}{\Lambda^{1/2}V_{r-a}(t)}_{L^2}.
\]
On a proposed finite maximal interval \(I=(0,T_*)\), the compatible finite
datum rectangles give the continuous \(H^1\) row traces
\[
 \norm{H^{(r)}}_{L^\infty(I)}
 \le\frac12\sum_{a=0}^r\binom ra
 \norm{V_a}_{C([0,T_*];H^1)}
 \norm{V_{r-a}}_{C([0,T_*];H^1)}<\infty.
\]
Every finite derivative of \(H\) belongs to \(L^2(I)\), since
\(T_*<\infty\). Thus \(H\in\bigcap_{r\in\Nzero}H^r(I)\), and
temporal Sobolev embedding gives its smoothness.  The identities below retain
the complete transport--pressure terms and expose, at every finite height, the
spatial energy carried by the same rectangle.

\begin{proposition}[Completed height recurrence in every temporal row]
\label{u:prop:completed-height-recurrence}
For every strict classical interval \((0,T)\), every \(r\in\Nzero\), and
every \(s\ge0\),
\begin{equation}
 \partial_tE_{r,s}=\cP_{r,s}-2\nu E_{r,s+1}.
 \label{u:eq:sobolev-height-balance}
\end{equation}
For every \(k\in\N\),
\begin{equation}
 \partial_t^kE_{r,1/2}
 =(-2\nu)^kE_{r,k+1/2}
 +\sum_{j=0}^{k-1}(-2\nu)^j
   \partial_t^{k-1-j}\cP_{r,j+1/2}.
 \label{u:eq:completed-height-every-temporal-row}
\end{equation}
In the base row \(r=0\),
\begin{equation}
 H^{(k)}
 =(-2\nu)^kE_{k+1/2}
 +\sum_{j=0}^{k-1}(-2\nu)^j
   \partial_t^{k-1-j}\cP_{j+1/2},
 \label{u:eq:completed-height-recurrence}
\end{equation}
and the highest spatial energy is the exact coordinate
\begin{equation}
 E_{k+1/2}
 =\frac{H^{(k)}-
   \displaystyle\sum_{j=0}^{k-1}(-2\nu)^j
   \partial_t^{k-1-j}\cP_{j+1/2}}
 {(-2\nu)^k}.
 \label{u:eq:completed-height-row}
\end{equation}
\end{proposition}

\begin{proof}
Pair the temporal recurrence for \(V_r\) with \(\Lambda^{2s}V_r\).
The pressure-complete nonlinear sum is exactly \(\cP_{r,s}\), while
\[
 \nu\operatorname{Re}\ip{\Lambda^sV_r}{\Lambda^s\Delta V_r}_{L^2}
 =-\nu\norm{\Lambda^{s+1}V_r}_2^2=-2\nu E_{r,s+1}.
\]
This proves \eqref{u:eq:sobolev-height-balance}.  Induction on \(k\) gives
the stronger identity
\begin{equation}
 \partial_t^kE_{r,s}
 =(-2\nu)^kE_{r,s+k}
 +\sum_{j=0}^{k-1}(-2\nu)^j
   \partial_t^{k-1-j}\cP_{r,s+j}.
 \label{u:eq:all-height-recurrence}
\end{equation}
Indeed, differentiation of the order-\(k\) identity followed by
\(\partial_tE_{r,s+k}=\cP_{r,s+k}-2\nu E_{r,s+k+1}\) produces the
order-\(k+1\) identity.  Taking \(s=1/2\) proves
\eqref{u:eq:completed-height-every-temporal-row}.  Taking \(r=0\) gives
\eqref{u:eq:completed-height-recurrence}; since \(\nu>0\), division by
\((-2\nu)^k\) gives \eqref{u:eq:completed-height-row}.
\end{proof}

\begin{definition}[Completed critical-height coordinates]
\label{u:def:completed-critical-height-coordinates}
Set
\begin{equation}
 R_0:=H=E_{1/2},
 \qquad
 R_k
 :=\frac{H^{(k)}-
   \displaystyle\sum_{j=0}^{k-1}(-2\nu)^j
   \partial_t^{k-1-j}\cP_{j+1/2}}
 {(-2\nu)^k}
 \quad(k\in\N).
 \label{u:eq:completed-critical-coordinate}
\end{equation}
Proposition~\ref{u:prop:completed-height-recurrence} gives the exact identity
\begin{equation}
 R_k=E_{k+1/2}
 =\frac12\norm{\Lambda^{k+1/2}u}_2^2
 \qquad(k\in\Nzero).
 \label{u:eq:completed-critical-coordinate-identity}
\end{equation}
\par
\end{definition}

\begin{definition}[Finite mixed rectangle]
\label{u:def:finite-rectangle}
Let \(0<T<T_*\) and \(N,M\in\Nzero\).  Restrict the whole-terminal record in
\eqref{u:eq:independent-datum-rectangle-record} to \((0,T)\):
\[
 \cR_{N,M}[u_0;T]
 :=\left(
 (V_n|_{(0,T)})_{n=0}^{N},
 (V_{n+1}|_{(0,T)})_{n=0}^{N};
 (Q_n|_{(0,T)})_{n=0}^{N}
 \right).
\]
Its associated velocity and pressure norms are
\begin{align}
 \mathfrak R_{N,M}(u;T)
 :={}&\sum_{n=0}^{N}
 \norm{V_n}_{L^2(0,T;H^{M+1})}^2
 +\sum_{n=0}^{N}
 \norm{V_{n+1}}_{L^2(0,T;H^{M-1})}^2,
 \label{u:eq:finite-rectangle-norm}\\
 \mathfrak Q_{N,M}(p;T)
 :={}&\sum_{n=0}^{N}
 \left(
 \norm{Q_n}_{L^1(0,T;H^M)}
 +\norm{\nabla Q_n}_{L^1(0,T;H^{M-1})}
 \right).
 \label{u:eq:finite-pressure-rectangle-norm}
\end{align}
The pressure-gradient coordinate is retained through
\eqref{u:eq:finite-pressure-rectangle-norm} and with every row constrained by
Proposition~\ref{u:prop:formal-mixed-recurrences}.
\end{definition}

\Needspace{14\baselineskip}
\begin{theorem}[Compatible finite rectangle norms]
\label{u:thm:datum-rectangles}

Let \(\nu>0\), \(u_0\in\HsSigma(\R^3)\), and let
\((T_*,u,p)\) be the maximal pressure-complete classical history generated
by \((u_0,\nu)\).  If \(T_*<\infty\), then, for every
\(N,M\in\Nzero\),
\begin{equation}
 \mathfrak R_{N,M}^*[u_0,\nu,T_*]
 :=\sup_{0<T<T_*}\mathfrak R_{N,M}(u;T)
 <\infty.
 \label{u:eq:datum-rectangle-bound}
\end{equation}
This assertion is coordinatewise in \((N,M)\); no bound uniform in all
derivative orders is asserted.  The quantity \(\mathfrak R_{N,M}^*\) is the
exact supremum for this history; for general data the theorem does not replace
it by a formula involving only finitely many norms of \(u_0\).
Equivalently, for every \(n,m\in\Nzero\),
\begin{equation}
 V_n\in L^2(0,T_*;H^{m+1}),
 \qquad
 V_{n+1}\in L^2(0,T_*;H^{m-1}).
 \label{u:eq:datum-adjacent-rows}
\end{equation}
Every coordinate in \eqref{u:eq:datum-adjacent-rows} is the corresponding
distributional derivative of the single maximal history.
\end{theorem}


\begin{proof}
Apply the compatible finite-rectangle statement of
Theorem~\ref{u:thm:datum-finite-blocks} on \(I=(0,T_*)\).  For fixed
\(N,M\), \eqref{u:eq:independent-adjacent-finite-rows} gives
\[
 V_n\in L^2(I;H^{M+1}),\qquad
 V_{n+1}\in L^2(I;H^{M+1})\hookrightarrow L^2(I;H^{M-1}),
 \qquad 0\le n\le N.
\]
Thus the finitely many selected coordinates have the finite total
\[
 C_{N,M}:=\sum_{n=0}^N
 \left(\norm{V_n}_{L^2(I;H^{M+1})}^2
       +\norm{V_{n+1}}_{L^2(I;H^{M-1})}^2\right)<\infty.
\]
Restriction decreases each norm, so
\[
 \mathfrak R_{N,M}(u;T)\le C_{N,M}
 \qquad(0<T<T_*).
\]
Taking the supremum proves \eqref{u:eq:datum-rectangle-bound}.
The number \(C_{N,M}\) is fixed before the cutoff \(T\) is chosen.
Its finiteness is coordinatewise and does not require a bound common to all
derivative orders.
\end{proof}

The displayed sum is the velocity norm of \(\cR_{N,M}\).  The next proposition
proves its equivalence with the adjacent-row memberships and identifies its
whole-interval value by monotone convergence.

\Needspace{14\baselineskip}
\begin{proposition}[Norm equivalence and full-interval realization]
\label{u:prop:datum-terminal-realization}

For the fixed maximal history, the family of bounds
\eqref{u:eq:datum-rectangle-bound} is equivalent to the adjacent-row statement
\eqref{u:eq:datum-adjacent-rows}.  Whenever either formulation holds, for every
\(N,M\in\Nzero\),
\begin{equation}
 \sum_{n=0}^{N}\left(
 \norm{V_n}_{L^2(0,T_*;H^{M+1})}^2
 +\norm{V_{n+1}}_{L^2(0,T_*;H^{M-1})}^2
 \right)
 =\lim_{T\uparrow T_*}\mathfrak R_{N,M}(u;T)
 =\mathfrak R_{N,M}^*[u_0,\nu,T_*].
 \label{u:eq:datum-terminal-monotone-realization}
\end{equation}
Every row is the corresponding distributional derivative of the single
pressure-normalized maximal history.
\end{proposition}

\begin{proof}
Assume first \eqref{u:eq:datum-adjacent-rows}.  For fixed \(N,M\), set
\[
 C_{N,M}:=
 \sum_{n=0}^{N}\left(
  \norm{V_n}_{L^2(0,T_*;H^{M+1})}^2
  +\norm{V_{n+1}}_{L^2(0,T_*;H^{M-1})}^2
 \right).
\]
The sum is finite because it contains finitely many instances of
\eqref{u:eq:datum-adjacent-rows}.  Restriction from \((0,T_*)\) to \((0,T)\)
gives
\[
 \mathfrak R_{N,M}(u;T)\le C_{N,M}
 \qquad(0<T<T_*),
\]
and hence \eqref{u:eq:datum-rectangle-bound}.

Conversely, assume \eqref{u:eq:datum-rectangle-bound}.  Every term in
\(\mathfrak R_{N,M}(u;T)\) is the integral of a nonnegative function.
Monotone exhaustion of \((0,T_*)\) gives
\eqref{u:eq:datum-terminal-monotone-realization}.  Taking \(N=n\) and \(M=m\)
yields \eqref{u:eq:datum-adjacent-rows} for each finite \((n,m)\).
Proposition~\ref{u:prop:formal-mixed-recurrences} identifies every entry with
the corresponding distributional derivative of the fixed pressure-normalized
history.
\end{proof}

\begin{corollary}[Critical-height exhaustion and Sobolev recovery]
\label{u:cor:datum-critical-height-completion}
For the datum-generated compatible finite rectangles of
Theorem~\ref{u:thm:datum-finite-blocks}, every completed critical coordinate
belongs to \(L^1(0,T_*)\), and
\begin{equation}
 \mathscr H_k[u_0,\nu,T_*]
 :=\norm{R_k}_{L^1(0,T_*)}
 =\frac12\norm{\Lambda^{k+1/2}u}_{L^2(0,T_*;L^2)}^2
 \le\frac12\mathfrak R_{0,k}^*[u_0,\nu,T_*]
 \qquad(k\in\Nzero).
 \label{u:eq:datum-completed-height-bound}
\end{equation}
For every \(m\in\Nzero\), the completed heights recover the spatial Sobolev
row through
\begin{equation}
 \norm u_{L^2(0,T_*;H^m)}^2
 \le
 2^m\left(
  T_*\norm{u_0}_2^2+2\norm{R_m}_{L^1(0,T_*)}
 \right)
 \le
 2^m\left(
  T_*\norm{u_0}_2^2+\mathfrak R_{0,m}^*[u_0,\nu,T_*]
 \right).
 \label{u:eq:datum-height-to-spatial-column}
\end{equation}
Thus the compatible critical-height exhaustion gives
\begin{equation}
 (R_k)_{k\in\Nzero}\in\prod_{k\in\Nzero}L^1(0,T_*),
 \qquad
 u\in\bigcap_{m\in\Nzero}L^2(0,T_*;H^m(\R^3)).
 \label{u:eq:datum-complete-height-and-spatial-columns}
\end{equation}
\end{corollary}

\begin{samepage}
\begin{proof}
Take the \(n=0,m=k\) finite rectangle in
\eqref{u:eq:independent-adjacent-finite-rows}; equivalently, take \(N=0\) and
\(M=k\) in \eqref{u:eq:datum-terminal-monotone-realization}.  Since
\(\abs\xi^{2k+1}\le\langle\xi\rangle^{2k+2}\),
\begin{align*}
 2\norm{R_k}_{L^1(0,T_*)}
 &=\norm{\Lambda^{k+1/2}u}_{L^2(0,T_*;L^2)}^2\\
 &\le\norm u_{L^2(0,T_*;H^{k+1})}^2
 \le\mathfrak R_{0,k}^*[u_0,\nu,T_*],
\end{align*}
which proves \eqref{u:eq:datum-completed-height-bound}.  For
\(\abs\xi\le1\),
\(\langle\xi\rangle^{2m}\le2^m\); for \(\abs\xi>1\),
\(\langle\xi\rangle^{2m}\le2^m\abs\xi^{2m+1}\).  Hence
\[
 \norm f_{H^m}^2
 \le2^m\left(\norm f_2^2+
 \norm{\Lambda^{m+1/2}f}_2^2\right).
\]
Integrating this inequality, using
\eqref{u:eq:energy-Linf} and
\eqref{u:eq:datum-completed-height-bound}, proves
\eqref{u:eq:datum-height-to-spatial-column}.  Since \(k\) and \(m\) are
arbitrary finite indices, \eqref{u:eq:datum-complete-height-and-spatial-columns}
follows.
\end{proof}
\end{samepage}

The preceding calculation recovers the same all-orders membership through the
completed heights.  For each finite \(k\), the quantity
\(R_k=E_{k+1/2}\) is computed from the velocity row of a sufficiently high
rectangle.  Compatibility under enlargement identifies all such evaluations
with the same history.  Taking every finite \(k\) gives
\eqref{u:eq:datum-complete-height-and-spatial-columns}.  Applying
Proposition~\ref{u:prop:spatial-column-reconstruction} to that complete spatial
row produces all temporal and pressure rows and hence
\eqref{u:eq:independent-mixed-intersection}.  Conversely, the
mixed-index exhaustion contains the velocity rows from which every finite
\(R_k\) is computed.  Both descriptions therefore use the same
datum-generated rectangle \(\cR_{N,M}\) at every finite pair \((N,M)\).

\begin{proposition}[Pressure completion]
\label{u:prop:rectangle-pressure-completion}
Under the hypotheses of Theorem~\ref{u:thm:datum-rectangles}, for every
\(N,M\in\Nzero\),
\begin{equation}
 \mathfrak Q_{N,M}^*[u_0,\nu,T_*]
 :=\sup_{0<T<T_*}\mathfrak Q_{N,M}(p;T)
 \le
 C_M^{\rm p}(2^{N+1}-1)\mathfrak R_{N,M}^*[u_0,\nu,T_*].
 \label{u:eq:whole-terminal-pressure-bound}
\end{equation}
\begin{samepage}
For every \(0\le n\le N\), the identities
\begin{align}
 Q_n
 &=R_iR_j\sum_{a=0}^n\binom naV_{a,i}V_{n-a,j},
 \label{u:eq:rectangle-pressure-identity}\\
 V_{n+1}
 &=\nu\Delta V_n
 -\sum_{a=0}^n\binom na(V_a\cdot\nabla)V_{n-a}-\nabla Q_n
 \label{u:eq:rectangle-velocity-identity}
\end{align}
hold respectively in
\(L^1(0,T_*;H^M)\) and \(L^1(0,T_*;H^{M-1})\), including when
\(M=0\).
\end{samepage}
\end{proposition}

\begin{proof}
For every \(0<T<T_*\),
Theorem~\ref{u:thm:formal-finite-row-pressure} gives
\[
 \norm{Q_n}_{L^1(0,T;H^M)}
 +\norm{\nabla Q_n}_{L^1(0,T;H^{M-1})}
 \le C_M^{\rm p}2^n\mathfrak R_{N,M}^*[u_0,\nu,T_*].
\]
Sum over \(0\le n\le N\), then let \(T\uparrow T_*\).  Monotone
convergence in the nonnegative norm integrals proves
\eqref{u:eq:whole-terminal-pressure-bound}.  Since
\(\nabla\cdot V_a=0\),
\[
 (V_a\cdot\nabla)V_{n-a}
 =\nabla\cdot(V_{n-a}\otimes V_a).
\]
Lemma~\ref{u:lem:sobolev-product} and Cauchy--Schwarz in time give
\[
 \norm{V_{n-a}\otimes V_a}_{L^1(0,T_*;H^M)}
 \le C_M^{\rm alg}
 \norm{V_a}_{L^2(0,T_*;H^{M+1})}
 \norm{V_{n-a}}_{L^2(0,T_*;H^{M+1})}.
\]
Thus every transport term belongs to
\(L^1(0,T_*;H^{M-1})\), also for \(M=0\).  Moreover,
\[
 \norm{\Delta V_n}_{L^1H^{M-1}}
 \le T_*^{1/2}\norm{V_n}_{L^2H^{M+1}}.
\]
Theorem~\ref{u:thm:datum-rectangles} places \(V_{n+1}\) in
\(L^2(0,T_*;H^{M-1})\), hence in \(L^1(0,T_*;H^{M-1})\).
The strict-interval identities are identities of distributions; their
full-interval representatives belong to the stated spaces.  This proves the
two identities on \((0,T_*)\).
\end{proof}

\begin{definition}[Finite rectangles and their restrictions]
\label{u:def:finite-window-space}
For \(N,M\in\Nzero\), the full-interval rectangle already defined in
\eqref{u:eq:independent-datum-rectangle-record} contains the rows
\begin{equation}
 \begin{aligned}
 V_n&\in L^2(0,T_*;H^{M+1}),\\
 V_{n+1}&\in L^2(0,T_*;H^{M-1}),\\
 Q_n&\in L^1(0,T_*;H^M),
 \end{aligned}
 \qquad 0\le n\le N.
 \label{u:eq:whole-terminal-rectangle-element}
\end{equation}
If \(N'\le N\) and \(M'\le M\), restricting to \((N',M')\) means discarding
the temporal coordinates with index larger than \(N'\) and applying
\[
 H^{M+1}\hookrightarrow H^{M'+1},
 \qquad H^{M-1}\hookrightarrow H^{M'-1},
 \qquad H^M\hookrightarrow H^{M'}
\]
to the retained coordinates.
\end{definition}

\begin{proposition}[Restriction compatibility]
\label{u:prop:rectangle-compatibility}
For \(N'\le N\) and \(M'\le M\),
\begin{equation}
 \left.\cR_{N,M}[u_0;T_*]\right|_{(N',M')}
 =\cR_{N',M'}[u_0;T_*].
 \label{u:eq:rectangle-compatibility}
\end{equation}
If
\((N'',M'')\le(N',M')\le(N,M)\), then
restricting first to \((N',M')\) and then to \((N'',M'')\) gives the same
rows as restricting directly to \((N'',M'')\).  Thus the datum-generated
rectangles form a compatible family.
\end{proposition}

\begin{proof}
All retained entries are the fixed distributions
\(\partial_t^nu\), \(\partial_t^{n+1}u\), and \(\partial_t^np\).
The Sobolev embeddings preserve those distributions.  Deletion of rows and
composition of canonical embeddings are associative.
\end{proof}
